% =========================================================
% PHẦN 1
% =========================================================
\section{Giới thiệu bài toán (Problem Statement)}

Chất lượng kết nối mạng có thể biến động theo thời gian do
lưu lượng truy cập, tình trạng đường truyền và các yếu tố
hạ tầng. Những biến động này ảnh hưởng đến độ trễ, mức độ
ổn định và trải nghiệm của người dùng.

Đề tài tập trung vào bài toán hồi quy (Regression) nhằm
dự đoán các chỉ số hiệu năng mạng như độ trễ (Latency),
độ dao động độ trễ (Jitter) hoặc thời gian phản hồi HTTP
dựa trên dữ liệu giám sát mạng.

Mục tiêu là xây dựng và đánh giá các mô hình Machine Learning
có khả năng ước lượng những chỉ số này, từ đó hỗ trợ giám sát
chất lượng kết nối và phát hiện sớm nguy cơ suy giảm hiệu năng.

Trong phạm vi yêu cầu của môn học, nhóm xây dựng mô hình
\textbf{hồi quy tuyến tính đơn biến}, sử dụng một biến đầu vào
$x$ để dự đoán một biến mục tiêu $y$:
\[
    \hat{y} = f_{\boldsymbol{\theta}}(x)
    = \theta_0 + \theta_1 x.
\]

Trong đó $\theta_0$ là hệ số chặn,
$\theta_1$ là hệ số góc; $x$ và $y$
đều có đơn vị mili giây (ms).

Thuật toán Gradient Descent được tự cài đặt bằng Python.
Các thư viện học máy cấp cao chỉ được sử dụng trong phần
so sánh tùy chọn.

\begin{itemize}
    \item Biến đầu vào $x$:
    \texttt{gateway\_rtt\_avg\_ms} --- độ trễ RTT trung bình
    đến Gateway, đơn vị mili giây (ms).
    \item Biến mục tiêu $y$:
    \texttt{rtt\_avg\_ms} --- độ trễ RTT trung bình đến
    máy chủ đích, đơn vị mili giây (ms).
    \item Lý do lựa chọn:
    cả hai đại lượng cùng phản ánh độ trễ mạng và có tương
    quan dương dự kiến. Độ trễ đến Gateway thường ổn định
    và dễ đo liên tục; nếu tồn tại quan hệ tuyến tính đủ
    mạnh, có thể dùng đại lượng dễ đo này để ước lượng độ
    trễ đến đích mà không phải gửi gói tin đến máy chủ từ
    xa trong mọi chu kỳ giám sát.
\end{itemize}
